% paper/sections/02b_surface_tension.tex
%
% §2.4  表面張力：界面条件とモデル概要
%   Extension PDE による圧力平滑化の詳細は第III部 §~\ref{sec:extension_pde} で述べる．
%   CSF モデルの完全な定式化は付録~\ref{app:csf_full} を参照．
%
% ═══════════════════════════════════════════════════════════════════════════════
\subsection{表面張力：界面条件とモデル概要}
\label{sec:csf}
% ═══════════════════════════════════════════════════════════════════════════════

表面張力は界面 $\Gamma$ 上の面力であり，NS 方程式（体積方程式）への組み込みには，
まず界面デルタ関数 $\delta_s$ を用いた特異体積力表現
（式~\eqref{eq:csf_intro}，§~\ref{sec:two_to_one}）：
\begin{equation*}
  \bm{f}_\sigma = -\sigma\kappa_{lg}\hat{\bm n}_{lg}\,\delta_s(\bx)
\end{equation*}
が数学的基盤となる．この連続表現には二つの読み方がある．
第一に，Brackbill \textit{et al.}~\cite{Brackbill1992} による
\textbf{連続表面力（CSF: Continuum Surface Force）モデル}では，
$\delta_s$ を幅 $\Ord{\varepsilon}$ の平滑化関数で置き換え，
一括一流体/Balanced--Force 経路の体積力
$\bm{f}_\sigma \approx \sigma\kappa_{lg}\bnabla\psi$ を構成する
（完全な定式化は付録~\ref{app:csf_full}）．
第二に，高密度比・低 $We$ を扱う主たる圧力ジャンプ経路では，同じ Young--Laplace 応力を
体積力として拡散させず，向き付き圧力ジャンプ
$j_{gl}=p_g-p_l=-\sigma\kappa_{lg}$ として分相 PPE に渡す
（§~\ref{sec:split_ppe_framework}）．
Ghost-Fluid Method（GFM）~\cite{Fedkiw1999, Kang2000} と
HFE（Hermite 場延長法；§~\ref{sec:field_extension}）は，
このジャンプを CCD/FCCD の片側面ジェットへ供給するための高次データ構成である．
したがって本稿で保持する階層は，
CSF が一流体/BF 経路で扱う拡散界面近似を与え，
分相圧力ジャンプ PPE が高密度比の圧力閉包を担う，というものである．
CSF 経路では $\Ord{\varepsilon^2} \approx \Ord{\Delta x^2}$ のモデル誤差が残るため
~\cite{Popinet2009, Francois2006}，
その精度上限は §~\ref{sec:pressure_accuracy_summary} で圧力ジャンプ型閉包と分けて整理する．

\noindent\textbf{$\delta_s \to \bnabla\psi$ の橋渡し．}
$\psi = H_\varepsilon(-\phi)$ より $\bnabla\psi = -H_\varepsilon'(-\phi)\,\bnabla\phi$ であり，
$H_\varepsilon'$ はコンパクト台を持つ平滑化 Dirac デルタ $\delta_\varepsilon(\phi)$ である
（$\int H_\varepsilon' = 1$，幅 $\Ord{\varepsilon}$；付録~\ref{app:csf_full}）．
Eikonal 条件 $|\bnabla\phi| = 1$ のもとで
$\delta_\varepsilon(\phi) \approx \delta_s$（coarea 公式）が成立するため，
$-\sigma\kappa_{lg}\hat{\bm n}_{lg}\,\delta_s \approx \sigma\kappa_{lg}\bnabla\psi$ という置き換えが正当化される．

\noindent\textbf{CSF 離散化の要点．}
CSF の体積力近似 $\bm{f}_\sigma = \sigma\kappa_{lg}\bnabla\psi$ は，
上記の $\bnabla\psi$ が界面近傍（中心遷移幅 $\sim 2\varepsilon$，95\% 遷移帯 $\sim 6\varepsilon$）に局在する性質を利用して，
表面張力を界面領域に自然に局在化する．
この $\bnabla\psi$ による局在化機構が，
Balanced--Force 条件（§~\ref{sec:balanced_force}）における
圧力勾配との演算子整合性の基盤となる．

\begin{tcolorbox}[defbox, title={CSF/Balanced--Force 経路へ渡す離散化上の注意}]
連続論では，静的平衡で圧力勾配と表面張力が釣り合う．
しかし離散化後に，評価位置，圧力 Poisson 方程式と速度補正の勾配，
発散と勾配の随伴性，面係数，界面ジャンプ補正のいずれかがずれると，
連続平衡は離散平衡として残らず，寄生流れの原因になる．
第2章ではこの事実を，CSF 体積力と Young--Laplace 圧力ジャンプが
同じ界面応力条件から来るという連続レベルの前提として固定する．
具体的な失敗モードの分解と設計原理は，
第~\ref{sec:balanced_force}節で扱う．
\end{tcolorbox}

一方，界面における圧力場は Young--Laplace 則に従う不連続を持つ．
界面圧力ジャンプ条件は：
\begin{equation}
  j_{gl} \equiv p_\text{gas} - p_\text{liq} = -\frac{\kappa_{lg}}{We}
  \label{eq:young_laplace_jump}
\end{equation}
（§~\ref{sec:notation} と同一の符号規約；無次元形式；$We$ の定義は §~\ref{sec:nondim}）．
この無次元形は有次元の Young--Laplace 則
$p_l^*-p_g^*=\sigma\kappa_{lg}^*$，すなわち
$j_{gl}^*=-\sigma\kappa_{lg}^*$ を
参照圧力 $\rho_l U^2$ で除し，$We = \rho_l U^2 L/\sigma$ と $\kappa_{lg} = L\kappa_{lg}^*$ を
代入して得られる（*-付き変数は有次元，*-無しは無次元；§~\ref{sec:nondim}）．

\noindent\textbf{表面エネルギーからの仮想仕事．}
Young--Laplace 則は，閉じた界面の表面エネルギー
\begin{equation}
  E_\Gamma[\Gamma]=\sigma |\Gamma|
  \label{eq:surface_energy_functional}
\end{equation}
の体積制約付き第一変分として読む．界面速度を $\bm w$，
外向き法線変位を $w_n=\bm w\cdot\hat{\bm n}$ とすると，
連続論では
\begin{equation}
  \delta E_\Gamma[\bm w]
  =
  \sigma\int_\Gamma \kappa\,w_n\,\mathrm{d}S,
  \qquad
  \delta V[\bm w]
  =
  \int_\Gamma w_n\,\mathrm{d}S .
  \label{eq:surface_energy_virtual_work}
\end{equation}
したがって静的平衡は，あるラグランジュ乗数 $\lambda$ に対して
\begin{equation}
  \delta(E_\Gamma-\lambda V)[\bm w]=0
  \quad\forall \bm w
  \qquad\Longleftrightarrow\qquad
  \sigma\kappa=\lambda
  \quad\text{on each closed component}
  \label{eq:surface_energy_constrained_equilibrium}
\end{equation}
である．ここで重要なのは，平衡判定が「真円か楕円か」という
形状名ではなく，\emph{体積を保つ全ての仮想変位に対して}
表面エネルギーの第一変分が消えるかで決まる点である．
非定曲率成分は，その形が楕円であるかどうかに依らず，
圧力反力で吸収できない毛管加速度として残る．
第~\ref{sec:split_ppe_capillary_range_projection}節では，
この連続仮想仕事を同じ面上の Riesz 表現と直交分解へ写し，
再初期化前の物理輸送終点上で実際の速度補正に用いる毛管補正量，成分体積反力，静的平衡判定を
圧力補正が使うものと同じ面計量で区別して定式化する．
この写像で保持しなければならない量は曲率そのものではなく，
仮想変位 $\bm w$ に対する仕事 $\delta E_\Gamma[\bm w]$ である．
離散化後には，許容される仮想変位は節点法線変位ではなく，
PPE 補正段と CLS 保存形移流が実際に使う面速度 $w_f$ である．
したがって，連続式~\eqref{eq:surface_energy_virtual_work} は
「界面上で曲率を評価して終わり」ではなく，
\[
  \delta E_\Gamma[T_f(q_T)w_f]
  =
  -\langle s_f,w_f\rangle_{M_A}
\]
という面空間の仕事恒等式へ変換してから補正段へ入れる．
ここで $T_f(q_T)$ は再初期化前の輸送終点における CLS 面輸送の線形化，
$M_A$ は圧力反力と同じ面計量である．
この段階で圧力仕事と異なる計量・異なる輸送終点・異なる発散を選ぶと，
連続 Young--Laplace 則が正しくても，離散静的平衡と非平衡駆動の判定は別の問題になってしまう．
壁付き計算では，この面空間はさらに no-slip 壁トレース制約
$C_ww_f=0$ を満たす許容空間へ制限される．
したがって壁反力は，圧力射影後に節点速度を上書きする補助処理ではなく，
表面エネルギー・圧力・体積制約の仮想仕事と同じ面速度状態空間を定義する
制約力である．

切断面圧力ジャンプでは，このジャンプを節点平均ではなく界面が通過する
切断面 $f\cap\Gamma_h$ で評価する．
面の両端値を $\psi^-,\psi^+$，$\kappa^-_{lg},\kappa^+_{lg}$ とすると
\begin{equation}
  \theta_f =
  \frac{1/2-\psi^-}{\psi^+-\psi^-},\qquad
  \kappa_{lg,f}=(1-\theta_f)\kappa^-_{lg}+\theta_f\kappa^+_{lg},
  \qquad
  B_{\Gamma,f}=\frac{j_{gl,f}}{H_f}
  =-\frac{\kappa_{lg,f}}{We\,H_f},
  \label{eq:cut_face_young_laplace_jump}
\end{equation}
である（$\theta_f$ は $[0,1]$ に制限し，$H_f$ は物理面距離）．
静止 no-slip 壁に界面端点が接する場合，端点の表面張力は壁反力で釣り合う制約力であり，
明示的な濡れ・接触角法則を導入しない限り領域内部の節点曲率を壁へ外挿してはならない．
したがって壁法線方向の閉包層では
\begin{equation}
  \partial_{\bm n_{\partial\Omega}}\kappa_{lg}=0
  \quad\text{（壁法線方向の毛管閉包層）}
  \label{eq:wall_capillary_curvature_closure}
\end{equation}
を用い，$\kappa_{lg,f}$ は領域内部側極限の界面/切断面幾何量として評価する．
これは曲率の上限制限や平滑化による安定化ではなく，no-slip 壁面トポロジ制約
式~\eqref{eq:wall_phase_topology_invariant} と同じ幾何制約から来る Young--Laplace 閉包である．

\noindent\textbf{毛管波時間スケール $\tau_\sigma$．}
表面張力が駆動する最短毛管波（波長 $\sim L$）の周期は，
古典的な毛管波の分散関係（DNS 文脈での CFL 取り扱いは~\cite{DennerVanWachem2015}）
\begin{equation}
  \omega^2 = \frac{\sigma k^3}{\rho_l+\rho_g}, \qquad k = \frac{2\pi}{L},
  \label{eq:capillary_dispersion}
\end{equation}
を $\tau_\sigma = 2\pi/\omega$ に代入して
\begin{equation}
  \tau_\sigma \equiv \sqrt{\frac{(\rho_l+\rho_g)\,L^3}{2\pi\sigma}}
  \label{eq:tau_sigma_def}
\end{equation}
で与えられる．
無次元化すると $\tau_\sigma/\tau_\mathrm{conv} = \Ord{\sqrt{We}}$ となり，
高表面張力（低 $We$）では $\tau_\sigma \ll \tau_\mathrm{conv}$ が時間刻みを律速する
（CFL 制約 $\Delta t \lesssim \tau_\sigma$；式~\eqref{eq:dt_sigma} §~\ref{sec:stability}）．

% ═══════════════════════════════════════════════════════════════════════════════
